\section{Correspondences and continuous model theory}
\label{sec:metric-logic}

How much of a statement about a metric space survives if all we have is a
finite table of distances? The diameter survives a small uniform distortion
of the table. So does the radius, even though its definition alternates two
quantifiers. An entire fragment of continuous first-order logic survives with
the \emph{same} error as its atomic formulas. This gives both a way to evaluate
sentences and a way to refute a proposed identification of two spaces. The
quantifiers are the real-valued $\sup$ and $\inf$ of metric model theory.
We use a relational, bounded fragment of the semantics of Ben Yaacov,
Berenstein, Henson and Usvyatsov.%
\footnote{I.~Ben Yaacov, A.~Berenstein, C.~W.~Henson and A.~Usvyatsov,
\emph{Model theory for metric structures}, especially section~3.
\href{https://math.univ-lyon1.fr/~begnac/articles/mtfms.pdf}{Authors' text}.}

\subsection{A category of finite observations}

Fix a finite signature of relation symbols with declared arities, including
a binary symbol for distance. A stage $A$ has vertices $0,\ldots,a$, with
$a\geq0$, and a rational $[0,1]$-valued table $P_A$ for each symbol $P$.
A metric stage also requires the distance axioms, which are finite tests.
The transfer argument needs only the tables, so it applies to pseudometrics.
A correspondence of grade $e\geq0$ from $A$ to $B$ is a Boolean matrix $R$
in which every row and column contains a one. For each $k$-ary $P$ and each
$k$-tuple of paired vertices $x_i R y_i$, require
\[
  |P_A(x_1,\ldots,x_k)-P_B(y_1,\ldots,y_k)|\leq e.             \tag{C}
\]
For a nullary relation this tests its two values. The finite test loops over
at most $(|A||B|)^k$ tuples for $P$. Multiplying all rational entries,
constants and grades by a common positive denominator converts the
comparisons in a given instance to integer arithmetic.

If $A\xrightarrow{R,e} B\xrightarrow{S,d} C$, put
$x(SR)z$ when $xRy$ and $ySz$ for some $y\in B$. Search up to $|B|$
for the middle vertex. Successive uses of row coverage and of column
coverage prove both cover conditions. For paired tuples in $SR$, choose
the first witnessing middle vertex in each coordinate. Apply (C) to each
side and the triangle inequality to get grade $e+d$. The diagonal has
grade zero. In a triple composite both bracketings assert the existence of
the same two successive intermediate vertices, so their Boolean matrices
are equal. The diagonal laws follow by substituting its single forced
coordinate. Thus correspondences compose in a category, with upper
bounds on distortion that add under composition. The same arrow can have
several upper bounds.

This gives an instance of the repository's
\emph{PresentedCategory} interface. Use natural numbers for validated
stage and matrix codes, extensional equality of Boolean entries for
parallel arrows, and the bounded matrix product above. Omit the grade
from the arrow code and carry (C) as a predicate. Every two-sided covering
relation has grade one in the bounded signature. The existing
presentation-wise Yoneda calculation applies to this instance once
its coding laws have been internalized. Formula transfer uses the
additional quantitative information in (C).%
\footnote{ProveIt, under \nolinkurl{Derivations/Lean/IPRADerivations/Foundations},
\texttt{PresentedCategory.lean} and \texttt{Yoneda.lean}. Repository paths in
this section are relative to \nolinkurl{Logic/PrimitiveRecursiveArithmetic}.}

\subsection{Quantifier transfer with a depth-independent error}

Formulas are built from predicates on variable tuples and rational
constants in $[0,1]$ by $1-u$, $u\mathbin{\dot-}q=\max(0,u-q)$,
$\min(u,v)$, $\max(u,v)$, $\sup_x u$, and $\inf_x u$. On a stage,
the last two operations take the maximum and minimum over its vertices.
Every stage is nonempty. A closed formula thus has an exact rational
stage value $v_A(\varphi)$. An open formula has a value for each assignment
of its free variables.

\paragraph{Transfer theorem.}
If $R$ has grade $e$ and the two assignments agree along $R$ coordinate
by coordinate, then the two values of any such $\varphi$ differ by at most
$e$. This holds for every fixed formula and arbitrary encoded stages.

\begin{proof}
An atom gives precisely (C). A constant gives zero discrepancy. Reflection
around $1/2$ and truncated subtraction are 1-Lipschitz, as are minimum
and maximum on pairs with the maximum norm. For the last assertion, if
$u_i\leq v_i+e$ for $i=1,2$, monotonicity bounds the minimum (respectively
maximum) of the two $u_i$ by the corresponding extremum of the two $v_i$
plus $e$. Interchange the lists for the other inequality.

Suppose $\varphi=\sup_x\psi$. Select a maximizing vertex $x$ of $A$ by a
finite scan, then select $y$ from its row of $R$. The inductive assertion
bounds $\psi_A(x)$ by $\psi_B(y)+e$, hence by $\max_B\psi_B+e$.
A maximizer of $B$ and coverage of its column give the converse bound.
For $\inf_x\psi$, a minimizer of $A$ and row coverage give
$\min_B\psi_B\leq\min_A\psi_A+e$, and a minimizer of $B$ gives the converse.
Every binder inherits the same $e$ as its body. Row coverage alone only
gives one direction of each quantified estimate.
\end{proof}

The proof is a derivation scheme by external recursion on the fixed
formula. For a binder, the internal loop invariant records an index in
the scanned prefix, its evaluated value, and the fact that this value
dominates (or is dominated by) the prefix. This is a bounded Boolean
test. With the formula fixed, the evaluator and the relation test are PR
by nested bounded scans and composition. Variables are represented by
de Bruijn indices in a numerical tuple. Adding a binder shifts the old
indices and places the tested vertex at index zero. A proof for a new
connective requires an explicit modulus. For example
$\min(1,u+v)$ propagates two errors to $2e$. The depth-independent
constant in the theorem uses the displayed fragment.

\paragraph{A uniform arithmetic theorem for encoded formulas.}
For this finite grammar the scheme also admits a single formula-code
version. Give the syntax in postorder, with $h$ instruction rows and
all child indices smaller than their parent. Choose $K$ one above the
largest variable index in any atom, or zero if none occurs. Use integer
units with a positive common denominator $D$, and reject malformed
instructions, arities and constants by bounded tests. For a stage of
size $a$, allocate $h a^K$ table cells with entries in $[0,D]$.
Column $u<a^K$ is an assignment in base $a$ with $K$ digits, treating
$a=1$ as the unique all-zero assignment. For a binder, prepend a
vertex to the assignment and retain the first $K$ coordinates.
Take the maximum or minimum over the $a$ resulting child entries.
When $K=0$ there is just the empty assignment. An atom indexes its
predicate table, and all other instructions use at most two child
rows. Thus there is a uniform PR evaluator on formula and structure
codes. Its evaluation tables can be packed in base $D+1$ with bound
$(D+1)^{h a^K}$. These bounds express finite table evaluation even
for signatures that include an arbitrarily large numerical table.

For related stages of sizes $a,b$, induct on the number of rows filled.
The invariant says that in every completed row, entries in related
columns differ by at most $e$. Checking it costs a scan over at most
$h a^K b^K$ triples and $K$ matrix entries per column test. In a binder
row the selected extremum and the row or column witness produce a pair
of extended related assignments; that child row has already passed the
invariant. Each remaining row rule is one of the numerical inequalities
above. The empty prefix is vacuous. This is a single QF induction on
a Boolean equation and gives the code-uniform transfer theorem, also
for a postorder circuit which shares subformulas. It interprets each predicate by lookup in its bounded numerical table.

\paragraph{Arithmetical meaning of a model value.}
Suppose a fixed PR family of stages $A_n$ comes with verified matrices
from $A_n$ to $A_{n+1}$ of grades at most $2^{-n-1}$. For a sentence
$\varphi$ compute $v_n=v_{A_n}(\varphi)$. For $m\geq n$ transfer and finite
summation give
\[
 |v_n-v_m|\leq\sum_{j=n}^{m-1}2^{-j-1}=2^{-n}-2^{-m}.
\]
The summation uses QF induction on a cross-multiplied rational inequality
with the endpoint $n$ fixed. Hence $v_n$ names a real value with error at
most $2^{-n}$, even when the stages have different denominator grids.
Its upper cut at a rational $q$ can be observed strictly whenever
$v_n+2^{-n}<q$. The analogous lower test uses $q<v_n-2^{-n}$.
The zero condition is the arithmetic assertion
$\forall n\ (v_n\leq 2^{-n})$, by the tail bound and nonnegativity.
It need not be decidable. Strict tests on opposite sides cannot both pass,
and any positive margin eventually appears in a test by the geometric
error bound. A second presentation related at accuracy tending effectively
to zero yields the same cuts by the transfer theorem and these bounds.

\paragraph{When this evaluates an ordinary metric structure.}
Suppose a nonempty compact metric structure $M$ is supplied with finite
$r_n$-nets of actual sample points, $r_n\to0$ with a stated PR rate,
and evaluations of its atomic predicates at those samples with error
$\eta_n\to0$ of stated PR rate. Assume each predicate has a specified
Lipschitz bound $L_P$ for the maximum metric on its tuples. The distance
predicate has $L_d=2$. Let $L$ be the largest bound. Relate a point of
$M$ to each sample within $r_n$ of it. A sample is related to itself and
coverage supplies a sample for each point. The atomic discrepancy is at
most $e_n=Lr_n+\eta_n$. With this relation the preceding induction works
for ordinary suprema as well. For the upper bound on a supremum, bound
\emph{every} value on $M$ by the finite maximum plus $e_n$. For the
reverse, use a maximizing sample and its copy in $M$. The analogous
bounds for an infimum reverse the comparisons. This proof needs no
assumption that a supremum is attained in $M$.

Here is a constructive reading if those suprema have yet to be defined.
Induct externally through the formula. At a quantifier, finite extrema on
two nets agree to within $e_n+e_m$. Pair samples through a common point
of $M$, use net coverage in both directions and the atomic bounds on
both sides, and apply finite formula transfer. Its proof also applies to
this externally specified relation, which supplies witnesses for every
row and column. For free variables retain their common assigned points.
These extrema have a regular Cauchy name after applying the supplied
rates. Its upper-bound property for a supremum follows by approximating
each argument with a net point to arbitrarily small error. Its leastness
follows because every selected maximizing sample is itself a point of
$M$, where the already interpreted body has value within $e_n$ of the
sample evaluation. The dual proof gives infima. This proves the semantic
bridge for the supplied presentation. Inside iPRA the statements concern
numerical stage comparisons and cuts for fixed presented names. To apply
the bridge, a classical description still requires nets, evaluation
algorithms and moduli.

\paragraph{Constructing the comparison matrices.}
There is a direct numerical construction when distances between sample
points of different levels can also be approximated. Suppose a PR table
$\widehat d_{nm}$ approximates all such distances within $\tau$, and
each sample of either level has a covering witness on the other level
with the advertised net radius. Write $r=\max(r_n,r_m)$. Retain the
pair $(i,j)$ if $\widehat d_{nm}(i,j)\leq r+\tau$. Every supplied
covering pair passes. Every retained pair has true distance at most
$r+2\tau$. Consequently (C) holds with
\[
       e=\eta_n+\eta_m+L(r+2\tau).
\]
Each assertion is a bounded scan or a direct instantiation of a supplied
metric or predicate estimate. Thus we obtain the required Boolean
matrix even when exact comparison of a real distance with $r$ is
undecidable. For example, thin the levels by their PR rates so that
$r_n,\eta_n$ at the new index $j$ are at most
$2^{-j-4}/(L+1)$, using prefix maxima of requested indices if needed.
Compute cross-distances at that same accuracy for successive levels.
The displayed grade is at most $2^{-j-1}$, as required by the Cauchy
argument. This supplies a numerical bridge for named spaces with common
distance evaluation.

\subsection{A sentence separating two continua of the same diameter}

Let $I=[0,1]$ and let $T=\mathbb R/2\mathbb Z$ carry its shorter-arc
metric. Both have diameter $1$. Consider
\[
                    \rho=\inf_x\sup_y d(x,y).
\]
For $m\geq1$ use units $1/(2m)$. The interval stage consists of
$0,\ldots,2m$ with ordinary distance in integer units. The circle stage
consists of $0,\ldots,4m-1$, with
$d(i,j)=\min(|i-j|,4m-|i-j|)$. Each point is within $1/(2m)$ of a
sample, and all sample distances are exact. For a constructive name
a rational approximation finer than one quarter of the mesh and rounding
it to a nearest sample give an explicit covering witness, including at a
circle seam by reduction modulo $4m$.
At an interval vertex $i$ the farther endpoint is at distance at least
$m$, while all vertices are within $m$ of the midpoint $m$.
At every circle vertex its antipode has distance $2m$, and no distance
exceeds $2m$. Thus the normalized interval and circle values are
$1/2$ and $1$ at every stage. The net bridge gives those same values
on the continua. Any proposed correspondence between these grids has
distortion at least $1/2$ even with arbitrarily fine mesh. The same
obstruction holds for a two-sided correspondence on the continua with
uniform distance distortion, by the supremum/infimum inequalities in
the transfer proof.

\subsection{Descent with a stated overlap error}

The repository's categorical methods paper proves exact gluing and
explicitly flags the missing error estimate for approximate overlap
data.%
\footnote{ProveIt, \nolinkurl{docs/proofs/foundations/ipra_categorical_methods/paper/sections/05_descent.tex},
subsection on a finite sheaf of function tables.}
Let a finite source stage $X$ be covered by masks $U_i$.
At each $x\in U_i$ there is a value $f_i(x)$ in a rational
pseudometric table $Y$. On overlaps assume
$d(f_i(x),f_j(x))\leq e$. Inside a patch, suppose a specified symmetric
nearness test $N(x,y)$ implies $d(f_i(x),f_i(y))\leq w$
when both points are in $U_i$. Finally require the finite cover
condition that every $N$-near pair lies in some common patch.

Choose the first patch $s(x)$ containing $x$, and set
$g(x)=f_{s(x)}(x)$. Every local value at $x$ is within $e$ of $g(x)$,
so every two patch selectors agree to within $e$. For any $N$-near
pair choose a common patch $i$. The three successive comparisons
\[
 g(x),\quad f_i(x),\quad f_i(y),\quad g(y)
\]
have costs $e,w,e$, so $d(g(x),g(y))\leq2e+w$. Any other global
table within $e$ of every patch is within $e$ of this particular selector
at each point, and two arbitrary such tables are within $2e$ of one
another. The $2e+w$ bound for arbitrary selectors is sharp with two singleton
patches followed by a two-point common patch. Take the common-patch
values $e,e+w$ in the rational line and the singleton values $0,2e+w$,
respectively. The first-cover selector takes both extreme values.
Rescale if a unit-diameter target is desired.

All hypotheses, including the common-patch condition, can be tested by
bounded table scans. Without the common-patch condition two nearby
points in disjoint singleton patches can be assigned unrelated values
even when $e=w=0$. A continuous application requires these estimates on
a refining family with accuracy tending to zero and coherence between
its selected global tables. For example, if matched vertices on two
levels admit corresponding patches whose local values agree within
$c$ in a common target, then their selected global values differ by at most
$e_n+c+e_{n+1}$, by the same three comparisons. With a summable PR
bound and PR tail modulus on these successive discrepancies one obtains
Cauchy names along specified compatible vertex names.

\begingroup\setlength{\emergencystretch}{3em}
\paragraph{Formalization.}
In \texttt{lean/Abstract.lean}, the correspondence laws underpin
\texttt{formula\_\allowbreak transport}. The calculations of both radii give
\texttt{circle\_\allowbreak line\_\allowbreak obstruction}, and the patch estimate is
\texttt{glue\_\allowbreak oscillation}. Its truth values and errors are integer units,
so arbitrary rational finite instances use one common denominator. The
file quantifies over Lean functions to state reusable lemmas. For a literal
object-theory derivation, instantiate the vertex sets by bounds, encode
each relation by a Boolean matrix and each predicate by its table,
totalizing malformed accesses by zero. The maximum/minimum loop needs
the extremum index and the bounded prefix domination test described
above. Relate variable assignments by a bounded conjunction over the
free positions. Derive the two quantifier inequalities by those extremum
indices and the two covering scans, then compile connective cases by
substitution. This produces an object proof for each fixed formula. The preceding postorder invariant gives the alternative route for one theorem over formula codes.
The Lean file checks the reusable semantic proof and the exact grid
calculation. The table-code compiler and its object-proof output for
this complete fragment remain to be mechanized.
\par\endgroup
